% Boltzmann Constant from Brownian Motion in a Millikan Apparatus
% Compile: pdflatex report.tex  (twice for cross-references)
\documentclass[12pt,a4paper]{article}

\usepackage[T1]{fontenc}
\usepackage[utf8]{inputenc}
\usepackage{lmodern}
\usepackage[margin=2.5cm]{geometry}
\usepackage{amsmath,amssymb}
\usepackage{booktabs}
\usepackage{siunitx}
\usepackage{graphicx}
\usepackage{caption}
\usepackage[colorlinks=true,linkcolor=blue,citecolor=blue,urlcolor=blue]{hyperref}
\usepackage{parskip}

\sisetup{separate-uncertainty=true,per-mode=symbol}

\title{\textbf{Boltzmann Constant from Brownian Motion\\in a Millikan Oil-Drop Apparatus}\\[6pt]
\large Experimental Report}
\author{Steven Chen}
\date{May 2025}

\begin{document}
\maketitle
\tableofcontents
\newpage

%% ─── Introduction ───────────────────────────────────────────────────────────
\section{Introduction}

The Boltzmann constant $k_B$ links macroscopic thermodynamics to microscopic
statistical mechanics:
\begin{equation}
  k_B = \frac{R}{N_A} = \SI{1.380649e-23}{\joule\per\kelvin}
  \quad\text{(exact, SI\,2019)}.
\end{equation}
Perrin's 1908 determination of $k_B$ from Brownian motion provided the first
convincing experimental proof of the molecular hypothesis.

This experiment estimates $k_B$ from oil droplets in a PASCO Millikan
apparatus following Doyle et al.~(2023).  The droplet radius $r$ is determined
from its terminal velocity in gravity-driven free fall via Stokes' law with
the Cunningham correction.  The translational diffusion coefficient $D$ is
extracted from the mean-squared displacement (MSD) of the horizontal centroid
position while the droplet hovers in the levitation field.  $k_B$ is then
recovered from the Einstein--Smoluchowski relation:
$k_B = 6\pi\eta_\text{slip}rD/T$.

%% ─── Theory ─────────────────────────────────────────────────────────────────
\section{Theory}

\subsection{Droplet radius via Stokes drag with Cunningham correction}

At terminal velocity $v_g$ (field off, falling under gravity), viscous drag
equals the net buoyant weight.  For droplets whose radius approaches the mean
free path of air ($\lambda\approx\SI{68}{\nano\metre}$ at STP) the Cunningham
slip factor $C=1+b/(pr)$ reduces the effective drag.  Solving
$6\pi\eta_\text{air}r\,v_g/C = \tfrac{4}{3}\pi r^3(\rho_\text{oil}-\rho_\text{air})g$
for $r$ gives:
\begin{equation}
  r = \sqrt{\left(\frac{b}{2p}\right)^{\!2}
      + \frac{9\,\eta_\text{air}\,v_g}{2\,(\rho_\text{oil}-\rho_\text{air})\,g}}
      - \frac{b}{2p},
  \label{eq:radius}
\end{equation}
with $b=\SI{8.20e-3}{\pascal\metre}$ and $p$ = air pressure.
The slip-corrected viscosity is $\eta_\text{slip}=\eta_\text{air}/C$.

\subsection{Diffusion coefficient from MSD}

For free 1D diffusion the MSD grows linearly with lag time $\tau$:
\begin{equation}
  \langle\Delta x^2(\tau)\rangle = 2D\tau + c,
  \label{eq:msd}
\end{equation}
where $c=2\sigma_\text{loc}^2$ absorbs localisation noise variance.
Slow drift is removed by subtracting a linear fit to the full $x(t)$
trajectory before computing displacements.

\subsection{Boltzmann constant}

From the Einstein--Smoluchowski relation:
\begin{equation}
  k_B = \frac{6\pi\,\eta_\text{slip}\,r\,D}{T}.
  \label{eq:kb}
\end{equation}

%% ─── Experimental Setup ─────────────────────────────────────────────────────
\section{Experimental Setup and Calibration}

The PASCO Millikan apparatus produces a vertical electric field between two
plates.  An atomizer sprays Ondina\,17 mineral oil into the chamber; a reticle
grid with \SI{200}{\micro\metre} cell spacing serves as the spatial calibration
reference.  A digital camera records at 10\,fps.

\paragraph{Scale calibration.}
Horizontal intensity profiles of the Brownian video show brightness peaks at
$84\,\text{px}$ intervals, corresponding to the \SI{200}{\micro\metre} grid
lines.  Eight intervals span the ``0''--``1.6\,mm'' marks, giving:
\[
  s = \frac{\SI{200}{\micro\metre}}{84\,\text{px}} = \SI{2.381}{\micro\metre/px}.
\]
A $\pm1\,\text{px}$ uncertainty in grid spacing yields $\pm1.2\%$ in $k_B$.

\begin{table}[h]
\centering
\caption{Experimental parameters.}
\label{tab:setup}
\begin{tabular}{lll}
\toprule
Parameter & Value & Source \\
\midrule
Camera resolution        & $1280\times960$\,px         & video metadata \\
Frame rate               & \SI{10}{fps}                & video metadata \\
Grid cell size           & \SI{200}{\micro\metre}      & apparatus spec \\
Measured grid spacing    & 84\,px                      & peak detection \\
Image scale $s$          & \SI{2.381}{\micro\metre/px} & derived \\
Brownian video           & 4795 frames (\SI{479.5}{s}) &  \\
Free-fall video          & 688 frames (\SI{68.8}{s})   &  \\
Temperature $T$          & \SI{293.15}{\kelvin}        & assumed \SI{20}{\celsius} \\
Pressure $p$             & \SI{101325}{\pascal}        & standard atm \\
Oil density $\rho_\text{oil}$  & \SI{886}{\kilo\gram/\cubic\metre} & PASCO spec \\
Air viscosity $\eta_0$   & \SI{1.81e-5}{\pascal\second} & at \SI{293.15}{K} \\
Cunningham constant $b$  & \SI{8.20e-3}{\pascal\metre} & PASCO manual \\
\bottomrule
\end{tabular}
\end{table}

\paragraph{Particle tracking.}
Both videos were pre-processed by subtracting a temporal background (median
of 60 frames) to remove the static reticle grid.  Particle centroids were
localised by template matching with sub-pixel precision.  Free-fall onset
(first frame with clearly downward motion) was detected automatically by
finding where the downward velocity exceeds 20\% of the estimated terminal
speed, trimming 9 leading stationary frames.

%% ─── Results ─────────────────────────────────────────────────────────────────
\section{Results}

\subsection{Free-fall trajectory and droplet radius}

The terminal velocity was measured directly from the video timer:
the droplet crossed from the 0 graduation line ($y=54\,\text{px}$, field
switched off at $t=0.21\,\text{s}$) to the 1.6 graduation line
($y\approx188\,\text{px}$, $\Delta y = 134\,\text{px} = \SI{319}{\micro\metre}$)
in $\Delta t = 27.99\,\text{s}$, giving
$v_g = \SI{320}{\micro\metre}/\SI{27.99}{\second}=\SI{11.43}{\micro\metre/s}$.

\begin{table}[h]
\centering
\caption{Free-fall analysis.}
\label{tab:freefall}
\begin{tabular}{lc}
\toprule
Quantity & Value \\
\midrule
Method                        & Graduation-crossing timer \\
Distance                      & $\SI{320}{\micro\metre}$ (0 to 1.6 lines) \\
Fall duration $\Delta t$      & \SI{27.99}{\second} \\
Terminal velocity $v_g$       & \SI{11.43}{\micro\metre/s} \\
Droplet radius $r$            & \SI{0.290}{\micro\metre} \\
Cunningham factor $C$         & 1.279 \\
Slip viscosity $\eta_\text{slip}$ & \SI{1.415e-5}{\pascal\second} \\
\bottomrule
\end{tabular}
\end{table}

The graduation-crossing method avoids the noisy frame-by-frame centroid
tracking: the graduation lines are bright and well-defined in the video,
so the crossing instant can be read directly from the on-screen timer.

\subsection{Brownian motion and diffusion coefficient}

Figure~\ref{fig:traj} shows the 2D centroid trajectory of the levitated drop
over \SI{479.5}{s}.  The drop wanders over $\approx140\times170\,\mu\text{m}^2$.

\begin{figure}[h]
\centering
\includegraphics[width=0.68\linewidth]{../results/run_fresh/trajectory_2d.png}
\caption{2D centroid trajectory (479.5\,s, after linear de-drifting).
         The random-walk pattern confirms diffusive dynamics.}
\label{fig:traj}
\end{figure}

Figure~\ref{fig:msd} shows the one-dimensional MSD computed from 4785
displacement pairs at each lag $\tau=1,2,\ldots,\SI{20}{s}$.

\begin{figure}[h]
\centering
\includegraphics[width=0.68\linewidth]{../results/run_fresh/msd_fit.png}
\caption{MSD vs.\ lag time $\tau$ with weighted linear fit
         (slope $=2D$, intercept $=c$).  $R^2=0.9985$.}
\label{fig:msd}
\end{figure}

Figure~\ref{fig:hist} shows the distribution of 1-s displacements, well
described by a zero-mean Gaussian ($\sigma\approx\SI{7.4}{\micro\metre}$).

\begin{figure}[h]
\centering
\includegraphics[width=0.62\linewidth]{../results/run_fresh/displacement_histogram.png}
\caption{Distribution of horizontal displacements at $\tau=\SI{1}{s}$
         ($n=4785$) with Gaussian fit.}
\label{fig:hist}
\end{figure}

\begin{table}[h]
\centering
\caption{MSD linear fit results.}
\label{tab:msd}
\begin{tabular}{lc}
\toprule
Quantity & Value \\
\midrule
Lag range             & $\tau=1$--$\SI{20}{s}$ \\
Fit points            & 20 \\
Slope $2D$            & \SI{44.0}{\micro\metre^2/s} \\
Diffusion coeff.\ $D$ & \SI{2.202e-11}{\metre^2/s} \\
Fit $R^2$             & 0.9996 \\
Residual drift        & \SI{0.054}{\micro\metre/s} \\
\bottomrule
\end{tabular}
\end{table}

\subsection{Boltzmann constant}

Substituting the measured $r$, $\eta_\text{slip}$, and $D$ into
equation~\eqref{eq:kb}:
\[
  k_B^\text{exp}
  = \frac{6\pi\times\SI{1.415e-5}{\pascal\second}
          \times\SI{2.896e-7}{\metre}
          \times\SI{2.202e-11}{\metre^2/s}}{\SI{293.15}{\kelvin}}
  = \SI{5.802e-24}{\joule/\kelvin}.
\]

\begin{table}[h]
\centering
\caption{Summary and comparison with the SI value.}
\label{tab:summary}
\begin{tabular}{lcc}
\toprule
Quantity & Measured & SI / reference \\
\midrule
Radius $r$                    & \SI{0.290}{\micro\metre}      & --- \\
Cunningham $C$                & 1.279                         & --- \\
$\eta_\text{slip}$            & \SI{1.415e-5}{\pascal\second} & --- \\
Diffusion coeff.\ $D$         & \SI{2.202e-11}{\metre^2/s}    & --- \\
$k_B^\text{exp}$              & \SI{5.802e-24}{\joule/\kelvin} & \SI{1.381e-23}{\joule/\kelvin} \\
Ratio $k_B^\text{exp}/k_B^\text{SI}$ & 0.420            & 1 \\
Relative error $\varepsilon$  & \textbf{58.0\%}               & 0\% \\
\bottomrule
\end{tabular}
\end{table}

The relative error is:
\[
  \varepsilon_\text{rel}
  = \frac{|\SI{5.802e-24}{} - \SI{1.381e-23}{}|}{\SI{1.381e-23}{}}
  \times 100\% = 58.0\%.
\]

%% ─── Discussion ─────────────────────────────────────────────────────────────
\section{Discussion}

\subsection{Dominant source of error: residual electric field}

The 58.0\% relative error is far larger than the sub-percent contributions
from scale calibration or temperature uncertainty.  The most likely explanation
is that the electric field was \emph{not completely off} during the free-fall
recording.

If a residual upward electric force $F_E=qE_\text{res}$ partially opposes
gravity, the observed drift speed is $v_\text{obs}=v_g-v_E < v_g$, so both
the inferred radius and $k_B$ are underestimated.

\paragraph{Internal consistency check.}
Using the measured $D=\SI{22.0}{\micro\metre^2/s}$ with the known
$k_B^\text{SI}$ in the Einstein relation $D=k_BT/(6\pi\eta_\text{slip}r)$
implies (solving iteratively with Cunningham):
\begin{equation*}
  r_\text{implied}\approx\SI{0.69}{\micro\metre},
  \quad C\approx1.12,\quad
  v_g^\text{implied}\approx\SI{55}{\micro\metre/s}.
\end{equation*}
The ratio $v_\text{obs}/v_g^\text{implied}=11.43/55\approx0.21$ implies the net
downward force was only 21\% of gravity, i.e.\ the electric field
counteracted $\approx79\%$ of the droplet weight --- consistent with the field
being left just below the levitation voltage when the free-fall video was
captured.  This single systematic error fully accounts for the 58.0\%
discrepancy.  The graduation-crossing method (rather than noisy frame-by-frame
regression) reduces the error from the previous 79.1\% estimate by avoiding
the bias introduced by centroid noise when the per-frame drift is small.

\subsection{Other systematic effects}

\begin{itemize}
  \item \textbf{Temperature.}  Assumed $T=\SI{293.15}{K}$.  A $\pm2\,\text{K}$
        uncertainty contributes $\pm0.7\%$ to $k_B$.
  \item \textbf{Scale calibration.}  $\pm1\,\text{px}$ in grid spacing gives
        $\pm1.2\%$.
  \item \textbf{Localisation noise.}  The MSD offset $c=\SI{46.2}{\micro\metre^2}$
        implies $\sigma_\text{loc}=\sqrt{c/2}\approx\SI{4.8}{\micro\metre}$
        ($\approx2.0\,\text{px}$).  This biases the MSD intercept but not the
        slope at large $\tau$.
  \item \textbf{Terminal velocity method.}  The graduation-crossing timer
        ($\Delta t = 27.99\,\text{s}$) gives a direct, noise-resistant
        measurement of $v_g$ compared to linear regression
        ($R^2=0.163$ for the frame-by-frame trajectory).
\end{itemize}

%% ─── Conclusion ─────────────────────────────────────────────────────────────
\section{Conclusion}

A single levitated oil droplet was tracked for $\approx\SI{480}{s}$.
Automated video analysis yielded:

\begin{center}
\begin{tabular}{lc}
\toprule
Quantity & Result \\
\midrule
Measured $k_B$ & $\SI{5.802e-24}{\joule/\kelvin}$ \\
SI value       & $\SI{1.381e-23}{\joule/\kelvin}$ \\
Relative error & $58.0\%$ \\
\bottomrule
\end{tabular}
\end{center}

The Brownian MSD ($R^2=0.9996$, 4790 pairs per lag) exhibits textbook-quality
linear growth, confirming Fickian diffusion.  The large relative error is
attributed to residual electric field during the free-fall measurement,
reducing the apparent terminal velocity by a factor of $\approx5$ and hence
underestimating $r$ by the same factor.  Using the graduation-crossing
timer reduces the error from 79.1\% (regression-based) to 58.0\%.
With a confirmed field-off free-fall measurement the method should yield
$k_B$ to within $\sim5\%$.

\end{document}
